% ISSUE 8 — Figure 3 Caption "DeFi AMM (Hard) = -78.04" Untraceable, Plus Impermanent-Loss Tier Clash
% Category: Open (confirmed known-bad figure/caption, fix source already identified)
% Source A: jmlr_paper_main.tex, Figure fig:r2_heatmap_clipped caption, l.1287-1300
% Source B: jmlr_paper_main.tex, tab:llm_ablation, l.1838-1901
% Source C: jmlr_paper_main.tex, Medium-tier definition (subsec:difficulty-classification), l.635-655
%   (full section also in file 07)

%% ---- SOURCE A: Figure 3 caption ----
\begin{figure}[htbp]
\centering
\includegraphics[width=\linewidth]{fig18_r2_heatmap_improved}
\caption{Mean far-extrapolation ($5\times$) $\Rsq$ by formula domain and
difficulty tier (Easy/Med/Hard), Core-15 benchmark. Dashes indicate no
equation exists for that domain--difficulty combination in this benchmark.
HypatiaX is not uniformly better than PySR-only: it improves on
Chemistry (Easy: $0.90$ vs.\ $-12.56$) but produces $-\infty$ --- a genuine
evaluation crash, not merely a poor fit --- on DeFi AMM (Hard) and Physics
(Hard), where PySR-only instead returns a finite (if poor, $-78.04$, or
perfect, $1.00$) value. The per-equation values underlying this aggregate
are given in Table~\ref{tab:llm_ablation}.}
\label{fig:r2_heatmap_clipped}
\end{figure}

%% ---- SOURCE B: tab:llm_ablation ----
Table~\ref{tab:llm_ablation} presents the full per-equation comparison between
PySR-only ($P$) and HypatiaX ($H$) on the Core-15 benchmark across three
extrapolation ranges (near $1.2\times$, canonical, far $5\times$).  The
Mann-Whitney U-test on far-$\Rsq$ vectors ($n=15$): $U=126$, $p=0.2948$
(two-sided); rank-biserial $r=-0.12$.  Run~B: $U=101.5$, $p=0.6833$.
The null result is stable across two hardware runs, indicating that the
HypatiaX advantage is domain-selective rather than universal.  Specifically:
\textbf{Physics} (Kinetic Energy, Gravitational Force, Ideal Gas Law) and
\textbf{DeFi Risk} (Value at Risk, Liquidation Price, Portfolio Variance) favour
HypatiaX; \textbf{Chemistry} (Arrhenius, Henderson-Hasselbalch) and
\textbf{DeFi AMM} (Impermanent Loss) favour PySR-only.  The null aggregate
$p$-value reflects this domain cancellation: gains on Physics/Risk are offset by
losses on Chemistry/AMM.  We report this honestly; the DeFi benchmark
($n=74$, §\ref{sec:results}) shows the effect size when evaluated on
HypatiaX's target domain.

\textbf{This table has been rebuilt from the underlying \texttt{\_merged.json}
data} (via \texttt{table6\_vs\_real\_data\_comparison.xlsx}), replacing the
withdrawn version. Six equations show \textit{nan} or $-\infty$ for HypatiaX
beyond the Train column --- these are missing or crashed evaluations, not
$\Rsq$ scores, and are reported as such rather than omitted. The originally-
published transcription (now withdrawn) had claimed finite values on all six
of these cells, including a fabricated $-12.5553$ ``identical collapse'' on
Arrhenius and a fabricated $-635$ on Michaelis-Menten, neither of which
exists in the underlying \texttt{\_merged.json} records.

\begin{table*}[htbp]
\centering
\caption{LLM Ablation: PySR Alone (P) vs.\ HypatiaX (H) on Core~15 ---
\textbf{rebuilt from \texttt{\_merged.json}} (real data, not the withdrawn
manuscript transcription). Extrap columns show $R^2$ at near ($1.2\times$),
medium (canonical), and far ($5\times$) out-of-distribution ranges.
\textit{nan}: evaluation never completed. $-\infty$: evaluation crashed.}
\label{tab:llm_ablation}
\small\setlength{\tabcolsep}{3pt}
\resizebox{\linewidth}{!}{%
\begin{tabular}{llcccccccc}
\toprule
\textbf{Equation} & \textbf{Domain} & \multicolumn{2}{c}{\textbf{Train $R^2$}} & \multicolumn{2}{c}{\textbf{Near $R^2$}} & \multicolumn{2}{c}{\textbf{Med $R^2$}} & \multicolumn{2}{c}{\textbf{Far $R^2$}} \\
\cmidrule(lr){3-4}\cmidrule(lr){5-6}\cmidrule(lr){7-8}\cmidrule(lr){9-10}
 & & P & H & P & H & P & H & P & H \\
\midrule
Allometric Scaling & Biology & 0.9977 & 0.9979 & 1.0000 & 0.9824 & 1.0000 & 0.9502 & 1.0000 & $-0.0644$ \\
Arrhenius & Chemistry & 0.9974 & 0.9979 & 0.9779 & 0.9977 & 0.4636 & 0.9957 & $-12.5552$ & \textbf{0.9028} \\
Constant Product & DeFi AMM & 0.9982 & 0.9985 & 0.9996 & $-1.807\times10^{4}$ & 1.0000 & $-5.416\times10^{34}$ & 0.9996 & $-1.852\times10^{46}$ \\
Gravitational Force & Physics & 0.9976 & 0.9889 & 0.9973 & $-\infty$ & 0.9998 & $-\infty$ & 0.9973 & $-\infty$ \\
Henderson-Hasselbalch & Chemistry & 0.9339 & 0.9986 & 0.8004 & \textit{nan} & 0.9850 & \textit{nan} & 0.8004 & \textit{nan} \\
Ideal Gas Law & Physics & 0.9976 & 0.9980 & 1.0000 & 0.9988 & 1.0000 & 0.9999 & 1.0000 & 0.9993 \\
Impermanent Loss & DeFi AMM & 0.9971 & 0.9977 & 0.7703 & 0.9639 & $-2.7832$ & $-80.4327$ & $-157.0776$ & $-8259.3$ \\
Kinetic Energy & Physics & 0.9968 & 0.9968 & 1.0000 & 1.0000 & 1.0000 & 1.0000 & 1.0000 & 1.0000 \\
Liquidation Price & DeFi Risk & 0.9974 & 0.9977 & 1.0000 & 0.9998 & 1.0000 & 0.9999 & 1.0000 & 0.9993 \\
Logistic Growth & Biology & 0.9975 & 0.9976 & 0.9997 & \textit{nan} & 0.9998 & \textit{nan} & 0.9996 & \textit{nan} \\
Michaelis-Menten & Biology & 0.9975 & 0.9979 & 0.9890 & \textit{nan} & 0.9951 & \textit{nan} & \textbf{0.6471} & \textit{nan} \\
Portfolio Std Dev & DeFi Risk & 0.9975 & 0.9975 & 1.0000 & 1.0000 & 1.0000 & 1.0000 & 1.0000 & 1.0000 \\
Price Impact & DeFi AMM & 0.9976 & 0.9976 & 1.0000 & 1.0000 & 1.0000 & 1.0000 & 1.0000 & 1.0000 \\
Rate Law & Chemistry & 0.9977 & 0.9981 & 1.0000 & \textit{nan} & 1.0000 & \textit{nan} & 1.0000 & \textit{nan} \\
Value at Risk & DeFi Risk & 0.9979 & 0.9979 & 0.9999 & 1.0000 & 1.0000 & 1.0000 & 0.9999 & 1.0000 \\
\bottomrule
\end{tabular}%
}% end resizebox
\par\vspace{4pt}
{\footnotesize\noindent
P = PySR-only; H = HypatiaX (PySR + LLM warm-start). \textbf{Bold}: cases
where the real data reverses the withdrawn table's headline claim (Arrhenius

%% ---- SOURCE C: Difficulty Classification (Easy/Medium/Hard defs) ----
\subsection{Difficulty Classification}\label{subsec:difficulty-classification}
\label{sec:bench_difficulty}

Tasks are categorised by difficulty into three tiers based on functional structure:

\begin{itemize}
  \item \textbf{Easy} ($n=24$): linear or simple rational relationships, such as
    Loan-to-Value, spot price from AMM reserve, collateral ratio, and simple staking APY.
    These tasks admit algebraic solutions with AST depth $\le 3$.

  \item \textbf{Medium} ($n=29$): nonlinear algebraic or exponential relationships,
    such as impermanent loss, Uniswap V3 liquidity range, constant-product price impact,
    and compounding staking returns.  These tasks require multi-step composition
    of algebraic operations.

  \item \textbf{Hard} ($n=21$): transcendental, probabilistic, or multivariate
    formulations.  This category includes Black--Scholes pricing (which requires
    the Gaussian CDF $\Phi$), all four Greeks, Portfolio Expected Shortfall for
    correlated assets, and multi-day VaR scaling.  These tasks are the primary
    source of LLM stochastic instability.
\end{itemize}
